% LoginMetricas — Informe de sustentación (ISO 25022 / 25023)
% Compilar en Overleaf: pdfLaTeX o XeLaTeX
% Autor: completar nombre, matrícula y fecha antes de entregar

\documentclass[11pt,a4paper]{article}

\usepackage[utf8]{inputenc}
\usepackage[T1]{fontenc}
\usepackage[spanish,es-tabla]{babel}
\usepackage{geometry}
\geometry{margin=2.5cm}
\usepackage{graphicx}
\graphicspath{{imagenes/}}
\usepackage{booktabs}
\usepackage{longtable}
\usepackage{array}
\usepackage{xcolor}
\usepackage{hyperref}
\usepackage{listings}
\usepackage{enumitem}
\usepackage{caption}
\usepackage{fancyhdr}

\hypersetup{
  colorlinks=true,
  linkcolor=blue!60!black,
  urlcolor=blue!70!black,
  citecolor=blue!60!black
}

\definecolor{codebg}{RGB}{245,245,250}
\lstset{
  backgroundcolor=\color{codebg},
  basicstyle=\ttfamily\small,
  breaklines=true,
  frame=single,
  rulecolor=\color{gray!40},
  keywordstyle=\color{blue!70!black},
  commentstyle=\color{gray},
  stringstyle=\color{teal!60!black},
  showstringspaces=false,
  tabsize=2,
  language=Python
}

\pagestyle{fancy}
\fancyhf{}
\lhead{LoginMetricas MVP}
\rhead{ISO/IEC 25022 y 25023}
\cfoot{\thepage}

% === Datos institucionales ===
\newcommand{\repoUrl}{https://github.com/OscarMeza24/LoginMetricas}
\newcommand{\ramaRepo}{master}
\newcommand{\fechaEntrega}{\today}

\newlength{\anchoCaptura}
\setlength{\anchoCaptura}{0.92\textwidth}
\newcommand{\incluirCalidad}[1]{%
  \begin{figure}[htbp]
    \centering
    \IfFileExists{imagenes/#1.png}{%
      \includegraphics[width=\anchoCaptura,keepaspectratio]{#1.png}%
    }{%
      \IfFileExists{imagenes/#1.jpg}{%
        \includegraphics[width=\anchoCaptura,keepaspectratio]{#1.jpg}%
      }{%
        \IfFileExists{imagenes/#1.jpeg}{%
          \includegraphics[width=\anchoCaptura,keepaspectratio]{#1.jpeg}%
        }{%
          \fcolorbox{black!35}{gray!8}{%
            \parbox[c][4cm][c]{0.85\textwidth}{\centering\small
              Sube \texttt{imagenes/#1.png} (o .jpg) a Overleaf}%
          }%
        }%
      }%
    }%
    \caption{#1}
  \end{figure}
}

\begin{document}

% --- Carátula ---
\begin{titlepage}
  \centering
  \vspace*{1.2cm}
  {\Large\bfseries UNIVERSIDAD LAICA ELOY ALFARO DE MANABÍ\par}
  \vspace{0.6cm}
  {\large\bfseries FACULTAD DE CIENCIA DE LA VIDA Y LA TECNOLOGÍA\par}
  \vspace{2.2cm}

  {\large\bfseries ASIGNATURA\par}
  \vspace{0.35cm}
  {\Large CALIDAD DE SOFTWARE\par}
  \vspace{2cm}

  {\large\bfseries TEMA\par}
  \vspace{0.35cm}
  {\Large\bfseries LoginMetricas\par}
  \vspace{0.25cm}
  {\normalsize MVP de autenticación y gestión de usuarios con FastAPI, GraphQL y React\par}
  \vspace{0.2cm}
  {\normalsize Cumplimiento ISO/IEC 25022:2016 e ISO/IEC 25023:2016\par}
  \vspace{2.2cm}

  {\large\bfseries AUTORES\par}
  \vspace{0.5cm}
  \begin{tabular}{c}
    MEZA MERA OSCAR ANDRÉS \\[0.35cm]
    FIGUEROA ARTEAGA LUIS EMILIO \\[0.35cm]
    CHILE SILVE CARLOS JUNIOR \\
  \end{tabular}
  \vfill

  {\large\bfseries DOCENTE\par}
  \vspace{0.35cm}
  ING. OSCAR BRIONES\par
  \vspace{1.2cm}

  {\large\bfseries PERIODO\par}
  \vspace{0.35cm}`'
  2026 -- 1\par
  \vspace{1cm}
  {\normalsize \fechaEntrega\par}
\end{titlepage}
\setcounter{page}{1}

\begin{abstract}
Este documento sustenta un \emph{Minimum Viable Product} (MVP) de autenticación y administración de usuarios, desarrollado con \textbf{FastAPI}, \textbf{GraphQL (Strawberry)} y \textbf{React (Vite)}. La evaluación de calidad sigue el marco \textbf{SQuaRE}: la norma \textbf{ISO/IEC 25023:2016} (calidad del \emph{producto}) y la \textbf{ISO/IEC 25022:2016} (calidad \emph{en uso}). Se documentan arquitectura, operaciones CRUD, medidas cuantitativas, evidencia en código, protocolo de demostración y resultados de pruebas automatizadas (\textbf{29/29} tests exitosos).
\end{abstract}

\tableofcontents
\newpage

%==============================================================================
\section{Introducción}
%==============================================================================

El proyecto \textbf{LoginMetricas} responde a la necesidad de un sistema seguro de inicio de sesión, registro, recuperación de contraseña y gestión de usuarios con roles (\texttt{admin} y \texttt{cliente}). El repositorio del equipo está disponible en:

\begin{center}
  \url{\repoUrl} (rama \texttt{\ramaRepo})
\end{center}

\subsection{Objetivos}
\begin{itemize}[noitemsep]
  \item Implementar autenticación con JWT y contraseñas hasheadas (bcrypt).
  \item Exponer un API GraphQL con CRUD completo sobre usuarios.
  \item Ofrecer interfaz web usable para demostrar calidad en uso (ISO 25022).
  \item Medir calidad de producto con métricas trazables y tests (ISO 25023).
\end{itemize}

\subsection{Stack tecnológico}
\begin{table}[h]
  \centering
  \caption{Componentes del sistema}
  \begin{tabular}{@{}lll@{}}
    \toprule
    \textbf{Capa} & \textbf{Tecnología} & \textbf{Ubicación} \\
    \midrule
    Interfaz & React 19 + Vite 6 & \texttt{web/src/pages/} \\
    Cliente API & Apollo Client & \texttt{web/src/services/} \\
    API & FastAPI + Strawberry GraphQL & \texttt{backend/app/resolvers.py} \\
    Negocio & Servicios Python & \texttt{backend/app/services.py} \\
    Seguridad & JWT + bcrypt & \texttt{auth.py}, \texttt{security.py} \\
    Persistencia & SQLite + SQLAlchemy & \texttt{models.py}, \texttt{database.py} \\
    \bottomrule
  \end{tabular}
\end{table}

\subsection{Relación entre ISO 25022 y 25023}
\begin{itemize}[noitemsep]
  \item \textbf{ISO/IEC 25023} evalúa propiedades del \emph{producto}: código, arquitectura, pruebas, seguridad implementada.
  \item \textbf{ISO/IEC 25022} evalúa la calidad cuando un \emph{usuario real} utiliza el software: efectividad, eficiencia, satisfacción, etc.
  \item Ambas normas se complementan dentro del marco SQuaRE (serie 25000).
\end{itemize}

%==============================================================================
\section{Arquitectura y flujo de datos}
%==============================================================================

El flujo de una petición autenticada es:

\[
  \text{UI (React)} \rightarrow \text{Apollo Client} \rightarrow \text{GraphQL} \rightarrow \text{Resolvers} \rightarrow \text{Services} \rightarrow \text{SQLite}
\]

En \texttt{main.py}, la función \texttt{get\_graphql\_context} inyecta la sesión de base de datos y la petición HTTP en cada operación GraphQL, lo que permite validar el token JWT y el rol antes de ejecutar mutaciones sensibles.

\subsection{Pantallas web (frontend)}
Login, registro, recuperación de contraseña, inicio, perfil, cambio de contraseña y gestión de usuarios (solo administrador). El acceso a rutas privadas se controla con \texttt{PrivateRoute} en \texttt{App.jsx}.

%==============================================================================
\section{Operaciones GraphQL y CRUD}
%==============================================================================

\begin{table}[h]
  \centering
  \caption{Mapeo CRUD $\rightarrow$ operaciones GraphQL}
  \small
  \begin{tabular}{@{}clll@{}}
    \toprule
    \textbf{CRUD} & \textbf{Operación} & \textbf{GraphQL} & \textbf{Autorización} \\
    \midrule
    C & Crear usuario & \texttt{register} & Público \\
    R & Leer uno & \texttt{getUser} & Autenticado \\
    R & Listar & \texttt{listUsers} & Solo \texttt{admin} \\
    U & Actualizar perfil & \texttt{updateUser} & Propio usuario o admin \\
    U & Cambiar contraseña & \texttt{changePassword} & Propio usuario o admin \\
    U & Reactivar & \texttt{activateUser} & Solo \texttt{admin} \\
    D & Soft delete & \texttt{deactivateUser} & Solo \texttt{admin} \\
    D & Hard delete & \texttt{deleteUser} & Solo \texttt{admin} \\
    \bottomrule
  \end{tabular}
\end{table}

Operaciones adicionales de autenticación: \texttt{login}, \texttt{verifyToken}, \texttt{requestPasswordReset}, \texttt{resetPassword}.

La autorización centralizada reside en \texttt{backend/app/auth.py} (\texttt{require\_admin}, \texttt{require\_self\_or\_admin}).

\subsection{Ejemplo de mutación (login)}
\begin{lstlisting}[language={},basicstyle=\ttfamily\footnotesize]
mutation {
  login(input: { email: "admin@admin.com", password: "Admin123!" }) {
    success
    message
    accessToken
    user { id email role }
  }
}
\end{lstlisting}

%==============================================================================
\section{Cumplimiento ISO/IEC 25023:2016 (Calidad del producto)}
%==============================================================================

La norma define medidas alineadas con ISO/IEC 25010. La Tabla~\ref{tab:25023} resume las medidas aplicadas al MVP.

\begin{longtable}{@{}p{2.2cm}p{2.4cm}p{3.2cm}p{1.6cm}p{1.4cm}p{2.6cm}@{}}
  \caption{Medidas de calidad de producto (ISO/IEC 25023)} \label{tab:25023} \\
  \toprule
  \textbf{Caract.} & \textbf{Subcar.} & \textbf{Medida} & \textbf{Objetivo} & \textbf{Resultado} & \textbf{Evidencia} \\
  \midrule
  \endfirsthead
  \multicolumn{6}{c}{\tablename\ \thetable\ -- \textit{Continúa}} \\
  \toprule
  \textbf{Caract.} & \textbf{Subcar.} & \textbf{Medida} & \textbf{Objetivo} & \textbf{Resultado} & \textbf{Evidencia} \\
  \midrule
  \endhead
  \bottomrule
  \endfoot
  Funcionalidad & Completitud & \% CRUD implementado & 100\% & \textbf{100\%} & register, get/list, update, delete \\
  Funcionalidad & Corrección & Tests unitarios OK & $\geq$95\% & \textbf{100\%} & 29 tests pytest \\
  Funcionalidad & Adecuación & Requisitos cubiertos & 100\% & \textbf{100\%} & Auth + roles + CRUD \\
  Confiabilidad & Disponibilidad & Health check & OK & \textbf{OK} & \texttt{/health} \\
  Confiabilidad & Tolerancia fallos & Rollback BD & Sí & \textbf{OK} & \texttt{services.py} \\
  Seguridad & Confidencialidad & bcrypt & 100\% & \textbf{100\%} & \texttt{security.py} \\
  Seguridad & Autenticidad & JWT Bearer & Sí & \textbf{OK} & \texttt{auth.py} \\
  Seguridad & Integridad & Validación entrada & 100\% & \textbf{100\%} & email, password, username \\
  Mantenibilidad & Modularidad & Capas separadas & $\geq$5 & \textbf{6} & models, services, resolvers... \\
  Mantenibilidad & Analizabilidad & Documentación & Completa & \textbf{OK} & README, ISO docs \\
  \bottomrule
\end{longtable}

\subsection{Métricas cuantitativas}

\textbf{Completitud funcional CRUD:}
\begin{equation}
  M_{\text{CRUD}} = \frac{\text{operaciones implementadas}}{\text{operaciones requeridas}} \times 100 = 100\%
\end{equation}

\textbf{Índice de correctitud (pruebas):}
\begin{equation}
  M_{\text{tests}} = \frac{\text{tests exitosos}}{\text{tests totales}} \times 100 = \frac{29}{29} \times 100 = 100\%
\end{equation}

\textbf{Validación de entrada:}
\begin{equation}
  M_{\text{valid}} = \frac{5}{5} \times 100 = 100\%
\end{equation}
(validaciones: email, contraseña fuerte, username único, email único, expiración JWT).

\textbf{Modularidad estructural:} $6/6$ módulos con responsabilidad única $= 100\%$.

\subsection{Evidencia en código (seguridad)}
\begin{lstlisting}
def require_admin(info, db):
    auth_user = require_auth(info, db)
    if auth_user.role != UserRole.ADMIN:
        raise GraphQLError("Se requiere rol de administrador")
\end{lstlisting}

\subsection{Pruebas automatizadas}
Comando de verificación:
\begin{lstlisting}[language=bash]
cd backend
source venv/bin/activate   # Python 3.11+
python -m pytest test_services.py -v
\end{lstlisting}
\textbf{Resultado documentado:} 29 pruebas pasando (100\%). Requiere Python 3.10 o superior.

%==============================================================================
\section{Cumplimiento ISO/IEC 25022:2016 (Calidad en uso)}
%==============================================================================

La norma define cinco características de calidad en uso. La Tabla~\ref{tab:25022} las relaciona con evidencia observable en la aplicación web.

\begin{table}[h]
  \centering
  \caption{Calidad en uso (ISO/IEC 25022)}
  \label{tab:25022}
  \small
  \begin{tabular}{@{}p{3cm}p{4.2cm}p{2cm}p{4.5cm}@{}}
    \toprule
    \textbf{Característica} & \textbf{Medida} & \textbf{Resultado} & \textbf{Evidencia en MVP} \\
    \midrule
    Efectividad & Tareas completadas sin ayuda & Alta & Login demo, perfil, CRUD admin en UI \\
    Eficiencia & Pasos para login & 2 campos + 1 clic & Botón ``Entrar con cuenta demo'' \\
    Satisfacción & Claridad de errores & Alta & Mensajes en español \\
    Satisfacción & Consistencia visual & Alta & DM Sans, paleta en \texttt{index.css} \\
    Libertad de riesgo & Seguridad percibida & Alta & Contraseña oculta, JWT, roles \\
    Libertad de riesgo & Privacidad & OK & Mensajes genéricos en login/reset \\
    Cobertura de contexto & Roles admin/cliente & OK & Pantalla Usuarios solo admin \\
    \bottomrule
  \end{tabular}
\end{table}

\subsection{Protocolo de demostración en clase}
\begin{enumerate}[noitemsep]
  \item Abrir \url{http://localhost:5173} y usar \textbf{Entrar con cuenta demo}.
  \item Editar perfil en \texttt{/profile} (efectividad + satisfacción).
  \item Gestionar usuarios: listar, desactivar, reactivar, eliminar (cobertura de contexto).
  \item (Opcional) Registrar usuario con contraseña débil $\rightarrow$ validación visible.
  \item (Opcional) GraphQL Playground en \url{http://localhost:8000/graphql}.
  \item Ejecutar \texttt{pytest} y mostrar 29 tests en verde (complemento 25023).
\end{enumerate}

\subsection{Credenciales de demostración}
\begin{center}
  \begin{tabular}{@{}ll@{}}
    \toprule
    Email & \texttt{admin@admin.com} \\
    Contraseña & \texttt{Admin123!} \\
    Rol & Administrador \\
    \bottomrule
  \end{tabular}
\end{center}
El usuario demo se crea o restablece al iniciar el backend (\texttt{seed\_default\_admin} en \texttt{database.py}).

%==============================================================================
\section{Procedimiento de ejecución}
%==============================================================================

\subsection{Backend}
\begin{lstlisting}[language=bash]
cd backend
python3.11 -m venv venv
source venv/bin/activate
pip install -r requirements.txt
python -m app.main
\end{lstlisting}
Verificar: \url{http://localhost:8000/health}

\subsection{Frontend web}
\begin{lstlisting}[language=bash]
cd web
npm install
npm run dev
\end{lstlisting}
Abrir: \url{http://localhost:5173} (proxy GraphQL en \texttt{/graphql}).

%==============================================================================
\section{Preguntas frecuentes del docente}
%==============================================================================

\begin{table}[h]
  \centering
  \small
  \begin{tabular}{@{}p{4.5cm}p{9.5cm}@{}}
    \toprule
    \textbf{Pregunta} & \textbf{Respuesta} \\
    \midrule
    ¿Por qué GraphQL? & Endpoint tipado; el cliente solicita solo los campos necesarios; adecuado para auth y CRUD. \\
    ¿Dónde está el CRUD? & \texttt{resolvers.py} (API) y \texttt{services.py} (reglas de negocio). \\
    ¿Cómo miden 25023? & Tabla de medidas + fórmulas + \textbf{pytest} reproducible. \\
    ¿Cómo miden 25022? & Protocolo de demo observable en navegador. \\
    ¿Soft vs hard delete? & \texttt{deactivateUser} $\rightarrow$ \texttt{is\_active=false}; \texttt{deleteUser} borra el registro. \\
    ¿Quién lista usuarios? & Solo rol \texttt{admin} (\texttt{require\_admin}). \\
    \bottomrule
  \end{tabular}
\end{table}

%==============================================================================
\section{Conclusiones}
%==============================================================================

El MVP \textbf{LoginMetricas} cumple el alcance funcional definido: autenticación JWT, registro, recuperación de contraseña, CRUD de usuarios con control por roles y documentación trazable. Las métricas ISO/IEC 25023 alcanzan \textbf{100\%} en completitud CRUD, correctitud (29/29 tests) y validación de entrada; las medidas ISO/IEC 25022 se demuestran mediante flujos de usuario verificables en la interfaz web. La evaluación SQuaRE queda respaldada por código, pruebas automatizadas y demostración en vivo.

%==============================================================================
\section*{Referencias}
\addcontentsline{toc}{section}{Referencias}
%==============================================================================

\begin{itemize}[noitemsep]
  \item ISO/IEC 25022:2016 --- Measurement of quality in use.
  \item ISO/IEC 25023:2016 --- Measurement of system and software product quality.
  \item ISO/IEC 25010:2011 --- System and software quality models (SQuaRE).
  \item Repositorio del proyecto: \url{\repoUrl}
\end{itemize}

\newpage
%==============================================================================
\section*{Anexos: capturas de pantalla}
\addcontentsline{toc}{section}{Anexos: capturas de pantalla}
%==============================================================================

Evidencia gráfica del MVP (carpeta \texttt{imagenes/} en Overleaf).

\incluirCalidad{calidad1}
\incluirCalidad{calidad2}
\incluirCalidad{calidad3}
\incluirCalidad{calidad4}
\incluirCalidad{calidad5}

\end{document}
